\documentclass[11pt,a4paper]{article}
\usepackage[T1]{fontenc}
\usepackage[utf8]{inputenc}
\usepackage{lmodern}
\usepackage{amsmath,amssymb,amsthm,mathtools,mathrsfs}
\usepackage{graphicx,float,xcolor,multicol}
\usepackage[shortlabels]{enumitem}
\usepackage[margin=27mm]{geometry}
\usepackage{microtype,needspace,placeins}
\usepackage{tikz,tikz-cd}
\usetikzlibrary{arrows.meta,calc,positioning,patterns,decorations.pathreplacing}
\usepackage[colorlinks=true,linkcolor=teal,urlcolor=teal,citecolor=teal]{hyperref}
\setlength{\parindent}{0pt}
\setlength{\parskip}{.6em}
\setcounter{secnumdepth}{0}
\allowdisplaybreaks
\raggedbottom
\providecommand{\cross}{\times}
\DeclareUnicodeCharacter{2020}{\ensuremath{\dagger}}
\DeclareMathOperator{\supp}{supp}
\DeclareMathOperator*{\esssup}{ess\,sup}
\providecommand{\chapter}{\section}
\newenvironment{tool}[2]{\subsection*{Tool #1 --- #2}}{}
\newcommand{\ToolRef}[1]{Tool #1}
\newcommand{\FigureTag}[2]{\textit{Figure #1}}
\newcommand{\Result}[1]{\subsection*{#1}}
\newcommand{\Application}[1]{\subsubsection*{Application\if\relax\detokenize{#1}\relax\else\space--- #1\fi}}
\newcommand{\ProofHeading}{\par\textit{Proof.}\space}
\newcommand{\ProofOf}[1]{\par\textit{Proof of #1.}\space}
\newcommand{\RemarkHeading}{\par\textbf{Remark.}\space}
\newcommand{\NoteHeading}{\par\textbf{Note.}\space}
\newcommand{\ExerciseHeading}{\par\textbf{Exercise.}\space}

\graphicspath{{figures/thermodynamic-formalism/}}
\begin{document}
\section*{Thermodynamical Formalism and Hausdorff Dimension}
% BLOG-CONTENT-BEGIN

How does thermodynamical formalism help us understand Hausdorff dimension?
Our motivation is to understand \cite{denker1991hausdorff} and to be
sufficiently equipped to extend its results to larger classes of functions
for which the conventional methods introduced by Bowen et al. fail.
\section{Overview of Ergodic Theory}
Given an object of a category with a self-morphism, we wish to study its
evolution. For a set $X$ with a group of automorphisms $G$, this is the realm
of group actions. For a compact metric space $X$ with a continuous
endomorphism $T$, it is the setting of topological dynamics. These objects
may be equipped with more refined structures. Finally, if $(X,\mu)$ is a
probability space and $T$ is measure-preserving, meaning that
$\mu(B)=\mu(T^{-1}(B))$, the evolution of $(X,\mu,T)$ belongs to ergodic
theory\footnote{For a complete treatment, see \cite{einsiedler2011ergodic}}.
One can think of $T$ as a deterministic link between the likelihoods of $B$
and $T^{-1}(B)$.

\paragraph{Examples.} 
\begin{enumerate}
    \item $(\mathbb{R}/\mathbb{Z}, \text{Lebesgue}, t\rightarrow t+\alpha)$
    \item $(\mathbb{R}/\mathbb{Z}, \text{Lebesgue}, t\rightarrow 2t)$
    \item $(\{0,1\}^{\mathbb{N}}, \prod_{\mathbb{N}}\mu_{\frac{1}{2}}, (x_0, x_1, \ldots) \rightarrow (x_1, x_2, \ldots))$ where $\mu_{\frac{1}{2}}(\{0\}) = \mu_{\frac{1}{2}}(\{1\}) = \frac{1}{2}$.
\end{enumerate}

We ask two fundamental questions about the evolution of a given measure-preserving system $(X, \mu, T)$.\\


1) What happens to $\{T^{\circ n}(x)\}_n$ or $\{T^{\circ n}(E)\}_n$ as $n\rightarrow \infty$?


2) What happens to $\{f\circ T^{\circ n}(E)\}_n$, for $f:X \rightarrow \mathbb{C}$, measurable, as $n\rightarrow \infty$? 

The answer to the first question leads us to the notions of recurrence and ergodicity. 
\subsection*{Theorem 2.1 — Poincare Recurrence} Given a measure-preserving system $(X, \mu, T)$ and a measurable set $B$ with $\mu(B) >0$,
$$\big|\{T^{\circ n}(x)\}_n\cap B\big| = \infty$$

for a.e. $x\in B$. Said plainly, the orbit of almost all $x\in B$ recurs in $B$ infinitely often.


Building on the theorem above, ask whether the same holds for $x\notin B$:
can $T$ make $x$ land in $B$? There is an immediate set-theoretic
obstruction. If $T^{-1}(B)=B$, then an orbit beginning outside $B$ cannot
enter it. This motivates the following definition.
\subsection*{Definition 2.1}
    A measure-preserving system $(X, \mu, T)$ is \textbf{ergodic} if $\forall B$, $T^{-1}(B) = B$, implies $\mu(B) \in \{0,1\}$. 


Ergodicity removes this obstruction, and recurrence then holds in the
following stronger sense:
$$\text{if }\mu(B)>0, \text{ then } \mu\Bigg(\bigcup_{n=1}^{\infty} T^{-n}(B)\Bigg) = 1$$


For a quantitative description, pass from the system to an operator on a Hilbert space using the following dictionary.
\begin{itemize}
    \item $L^2(\mu) \longleftrightarrow (X,\mu)$ (replace $X$ with a naturally occurring Hilbert space).
    \item $U_T \longleftrightarrow T$ (erect an operator $U_T$ on that Hilbert space).
    \item $\big(U_T:L^2(\mu) \rightarrow L^2(\mu)\big)\longleftrightarrow \big(T:X\rightarrow X\big)$ (study the action of $T$ on $X$ through the action of $U_T$ on $L^2(\mu)$).
\end{itemize}

Here, $U_T(f):= f\circ T$ is called the Koopman operator. For
$f_1,f_2\in L^2(\mu)$,
$$\langle U_Tf_1, U_Tf_2\rangle := \int_X (f_1\circ T)\overline{(f_2\circ T)}d\mu = \int_Xf_1\overline{f_2}d\mu$$

Since $T$ is measure-preserving, this implies
$$\langle f_1, (U_T^*U_T-\text{Id})f_2\rangle = 0$$

which forces, 
$$\sup_{||f_1||, ||f_2||\leq 1} \bigg(\big|\langle f_1, (U_T^*U_T-\text{Id})f_2\rangle\big|\bigg) = ||U_T^*U_T-\text{Id}|| = 0$$

Therefore, $U_T$ is unitary. The following theorem describes ergodicity from an operator theoretic perspective. 
\subsection*{Theorem 2.2}
    $(X, \mu, T)$ is ergodic if and only if $U_T(f) = f$ implies $f$ is some constant function $c$.


The answers to the second question appear as ergodic theorems. The common feature of these theorems is that they all study the quantity,
$$A_N^f:= \frac{1}{N}\sum_{n=0}^{N-1}U_T^n(f)$$

as $N \rightarrow \infty$. The principal statements are as follows.
\subsection*{Theorem 2.3 — von-Neumann Ergodic Theorem} Let the following be a projection operator,
$$P_T:L^2(\mu) \rightarrow \{g\in L^2(\mu)  | U_Tg= g\}$$

then, $\forall f\in L^2(\mu)$, 
$$A_N^f \rightarrow P_T(f)$$

in the $L^2$-norm.

\subsection*{Theorem 2.4 — Mean Ergodic Theorem} Given $f \in L^1(\mu)$, we have,
$$A_N^f \rightarrow f' (\in L^1(\mu))$$

in the $L^1$-norm, where $U_Tf' = f'$. 

As we have the following containment,
$$L^{\infty}(\mu) \subseteq \cdots \subseteq L^2(\mu) \subseteq L^1(\mu)$$

we see that Theorem $2.4$ extends Theorem $2.3$. 
\subsection*{Theorem 2.5 — Birkhoff's Ergodic Theorem} If $f \in L^1(\mu)$, then,
$$A_N^f(x) \rightarrow f^*(x)$$

pointwise. Furthermore, $f^* \in L^1(\mu)$, $U_Tf^* = f^*$, and,
$$\int_Xf^*d\mu = \int_Xfd\mu$$

\textbf{Remark 2.1.} If $(X,\mu, T)$ is ergodic, then $f^* = c$, and this implies, 
$$\int_Xf^*d\mu = c = \int_Xfd\mu$$

Therefore, 
$$\lim_{n\rightarrow\infty} \bigg(\frac{1}{n}\sum_{i=0}^{n-1}f(T^i(x))\bigg) = \int_Xfd\mu$$

If $f=\chi_A$, then $A_N^f$ is the frequency with which the orbit of $x$
visits $A$, called the time average. The right-hand side is the space average.
Birkhoff's theorem therefore exhibits a duality between time and space
averages.



In summary, ergodic systems separate events on average, weakly mixing systems
separate events asymptotically outside a density-zero set, and mixing systems
separate events asymptotically. This gives the containment
$$\text{Mixing } \subset \text{ Weakly Mixing } \subset \text{ Ergodic}$$


Let us now apply the tools developed thus far to understand topological dynamical systems. Recall that this setting consists of a compact metric space $(X, d)$, and a continuous self-map $T:X\rightarrow X$. To proceed one must induce the structure of a measure-preserving system on $(X, d, T)$. The Riesz representation theorem gives the bijective correspondence
$$\big(\Lambda \in C(X)^*\big) \longleftrightarrow \Bigg(f\rightarrow \int_Xfd\mu\Bigg)$$

Let $M(X)$ be the space of probability measures on $X$, which by the correspondence above can be identified with functionals on $C(X)$ with $*$-topology. Observe the following,
\begin{enumerate}
    \item Given a sequence $\{\mu_n\}_n$ from $M(X)$, with $\mu_n\rightarrow \mu$, we have, 
    $$\int_Xfd\mu_n \rightarrow \int_Xfd\mu$$
    
    If $f \equiv 1$, then $\mu(X) = 1$, and similarly one can also prove that $\mu > 0$. Therefore, $M(X)$ is closed.
    \item As $X$ is compact, the following holds,
    $$\Bigg|\int_xfd\mu\Bigg|\leq ||f||_{\infty}\int_X1d\mu$$
    
    which implies that $||\mu|| \leq 1$, giving us that $M(X)$ is bounded.
    \item By Banach--Alaoglu theorem, $M(X)$ is compact.
\end{enumerate}


The set 
$$M^T(X) := \{\mu\in M(X) |  \mu \text{ is }T-\text{invariant}\}$$

also satisfies $(1), (2), $ and $(3)$. These properties yield the following theorems.
\subsection*{Theorem 2.6}
    $M^T(X) \neq \emptyset$. That is, one can do ergodic theory on $X$

\subsection*{Theorem 2.7}
    Ergodic elements of $M^T(X)$ are its extreme points (call it $E^T(X))$.

\subsection*{Theorem 2.8 — Ergodic Decomposition} Given $(X, d, T)$, $\mu\in M^T(X)$, there exists unique $\lambda$ on $E^T(X)$ such that,
\begin{itemize}
    \item $\lambda(E^T(X)) = 1$
    \item The following convex linear integral representation holds,
    $$\int_Xfd\mu = \int_{E^T(X)}\bigg(\int_Xfd\nu\bigg)d\lambda$$
\end{itemize}

This has two consequences for the topological dynamical system.
\subsection*{Theorem 2.9 — Unique Ergodicity}
$A_N^f(x) \rightarrow C_f$

is independent of $x$ ($\forall f \in C(X)$) iff $|M^T(X)| =1$.

The second consequence is a notion of uniform distribution for sequences in
$X$.
\subsection*{Definition 2.4}
A sequence $\{x_n\}_n \subset X$ is said to be \textbf{equidistributed} in $(X, d, T, \mu)$, if $\forall f\in C(X)$, 
$$\lim_{n\rightarrow \infty} \frac{1}{n}\sum_{i=1}^{n}f(x_i) = \int_Xfd\mu$$

Further, $x\in X$ is said to be \textbf{generic} if $\{T^{\circ n}(x)\}_n$ is equidistributed.

\subsection*{Theorem 2.10}
    $\mu\in M^T(X)$ is ergodic implies that almost every $x\in X$ is generic. 


In this overview, ergodic theory arose from the evolution of a
measure-preserving system. Questions about its orbits led to recurrence and
ergodic theorems, which in turn gave a precise notion of equidistribution for
topological dynamical systems.

\section{Overview of Entropy}
Suppose events $A_1,A_2,\ldots,A_n$ have probabilities
$p_1,p_2,\ldots,p_n$. If $p_1=99\%$ and the remaining $1\%$ is distributed
among $\{p_i\}_{i>1}$, there is little uncertainty. If instead $p_i=1/n$ for
every $i$, the uncertainty is much greater. How can this be quantified? The
natural axioms for a function $H(p_1,p_2,\ldots,p_n)$ are the following.

\begin{enumerate}
    \item If $f(N):= H(1/N, 1/N, \ldots, 1/N)$, then $f$ is decreasing.
    \item $f(ML) = f(M) + f(L)$
    \item $H(p_1, \ldots, p_n) = h(p_A, p_B) + p_AH(p_1/p_A, \ldots, p_r/p_A) + p_BH(p_{r+1}/p_B, \ldots, p_N/p_B)$, where $A$ is the collection of the events $A_1, \ldots, A_r$, and $B$ is the collection of the remaining events.
    \item The function $p\rightarrow H(p, 1-p)$ must be continuous on $(0,1)$.
\end{enumerate}

The axioms above roughly capture our intuition on uncertainty: As $N$ increases, the fact that we demand $f(N)$ to decrease captures the fact that each outcome has almost no chance of occurring, hence inviting no dynamic expectation regarding the outcome. Axiom $2$ talks about a scenario in which there are $M$, $L$ events (each equally likely), and we are to choose one from the set of $M$ events, and the other from the set of $L$. In all, there are $ML$ equally likely choices. We demand the complexities of the two choices, captured by $f(M)$ and $f(L)$, to add up to $f(ML)$, to indicate that one choice does not affect the other in any quantifiable way; they simply add up. Axiom $3$ is just a strengthening of the scenario that Axiom $2$ addresses, namely, suppose the first choice is made with probability $p_A$, and the second with probability $p_B$, and further that the events are not equally likely, then we demand the uncertainties to add up like above. The rationalization for Axiom $3$ can be extrapolated from that of Axiom $2$. Finally, Axiom $4$ is a regularity demand which captures the nature of the uncertainties not having strange jumps (or discontinuities). The upshot of these Axioms is the following theorem.
\subsection*{Theorem 3.1}
    The functions satisfying Axioms $1, 2, 3, 4$ are all of the form,
    $$H(p_1, p_2, \ldots, p_n) = -c\sum_{i=1}^np_i\log{p_i}$$


One can use this to inform our choices for a general probability space $(X,\mu)$. 
\subsection*{Definition 3.1}Let $\alpha = \{A_1, A_2\ldots\}$ be a finite or countable partition of $X$. The \textbf{information function} of $\alpha$ is defined to be,
$$I(\alpha)(x):= -\sum_{A\in \alpha}\chi_A(x)\log{\mu(A)}$$  

\subsection*{Definition 3.2}
    With $\alpha$ as above, the \textbf{entropy} of $\alpha$ is defined to be,
    $$H(\alpha):= -\int I(\alpha)d\mu = -\sum_{A\in \alpha}\mu(A)\log{\mu(A)}$$


The information function measures how much the partition $\alpha$ narrows
down the location of $x\in X$. Large partition elements convey less
information, which is why $\chi_A(x)$ is weighted by $\log(1/\mu(A))$.
Entropy is the average value of this information function. Now suppose that
$\mathcal A\subseteq\mathcal B$ is a $\sigma$-subalgebra. For any
$f\in L^1(\mu)$, we have
$$\Bigg(\int_Xfd\big(\mu|_{\mathcal{A}}\big)\Bigg) \ll \mu$$

Let $E(f|\mathcal{A})$ be its Radon-Nikodym derivative. We now have an operator, 
$$E(\cdot|\mathcal{A}):L^1(\mu) \rightarrow L^1(\mu|_{\mathcal{A}})$$

such that, $\forall A\in \mathcal{A}$, and $\forall f\in L^1(\mu)$
$$\int_AE(f|\mathcal{A})d\mu|_{\mathcal{A}} = \int_Afd\mu$$

\subsection*{Definition 3.3}
    The \textbf{conditional probability} of $B\in \mathcal{B}$ given $\mathcal{A}$ is,
    $$\mu(B|\mathcal{A}) := E(\chi_B|\mathcal{A})$$

\subsection*{Definition 3.4}
    The \textbf{conditional information function} of $\alpha$ given $\mathcal{A}$ is,
    $$I(\alpha|\mathcal{A})(x):= -\sum_{A\in \alpha}\chi_A(x)\log{\mu(A|\mathcal{A})}$$
    
    and the \textbf{conditional entropy} of $\alpha$ is defined to be,
    $$H(\alpha|\mathcal{A}):= -\int I(\alpha|\mathcal{A})d\mu = -\sum_{A\in \alpha}\mu(A|\mathcal{A})\log{\mu(A|\mathcal{A})}$$


Now, let us upgrade our situation to a measure-preserving system $(X, \mu, T)$. Keeping in mind the terms introduced, we want to capture the complexity of the evolution induced by $T$ in $(X, \mu)$. Now, the information given by $\alpha$ to locate $x\in X$ is not absolute, for one can use $T$. Denote by $T^{-1}\alpha$ the partition $\{T^{-1}(A)\}_{A\in \alpha}$. For partitions $\alpha_1, \alpha_2$, denote by $\alpha_1\vee\alpha_2$ the partition $\{A_1\cap A_2\}_{A_i\in \alpha_i}$. Write,
$$H_n(\alpha) := H\Bigg(\bigvee_{i=0}^{n-1}T^{-1}\alpha\Bigg)$$

This measures the entropy accumulated through the first $n$ uses of $T$.
Successive uses of $T$ are not independent from the viewpoint of entropy:
in general $H_{n+m}(\alpha)\neq H_n(\alpha)+H_m(\alpha)$. Instead, the
sequence $\{H_n(\alpha)\}_n$ is subadditive:
$$H_{n+m}(\alpha)\leq H_{n}(\alpha) + H_{m}(\alpha)$$ 




The entropy of a probability space measures uncertainty, or the average
complexity of the information function. We now ask whether the entropy of a
measure-preserving system is likewise associated with a potential. This leads
to the following notion.
\subsection*{Definition 3.7}
    Let $(X, \mu)$ and $(X', \mu')$ be probability spaces. Let $T:X\rightarrow X'$ be a measurable map. We say a function $J$ is a \textbf{weak Jacobian} if there exists a measure $0$ set $E$ such that for all $A \subset X\backslash E$, on which $T$ is injective, the set $T(A)$ is measurable and,
    $$\mu(T(A)) = \int_AJd\mu$$
    
    We say, $J$ is a \textbf{Jacobian}, if the above holds for all $A\subset X$.

\subsection*{Theorem 3.6 — Rokhlin's Entropy Formula} Let $(X, \mu, T)$ be a measure-preserving system and a Lebesgue space, where $T$ is essentially countable to one\footnote{for details on the hypotheses refer \cite{przytycki2010conformal}}. Suppose that on each component $A$ of the ergodic decomposition, the restriction $T_A$ has a countable one-sided generator of finite entropy. Then,
$$h_{\mu}(T) = \int \log{J_{\mu}}d\mu$$


In short, the entropy of a measure-preserving system $(X,\mu,T)$ captures the
average logarithmic rate of expansion of $T$. The analogous notion for a
topological dynamical system is topological entropy; see
\cite{viana2016foundations}. For a fuller treatment of the material here, see
\cite{przytycki2010conformal}.

\section{Gibbs Measure}
This section gives an overview of Gibbs measures and their relation to the pressure functional; see \cite{przytycki2010conformal} for details. The setting of our story is the following:

\begin{enumerate}
    \item $X$, a compact metric space (with metric $\rho$)
    \item $T$, an open continuous distance expanding map (topologically transitive)
    \item $\varphi\in C(X)$ is H\"older-continuous.
\end{enumerate}
We begin in this generality and make three observations about the setting.

First, note that $T:X \rightarrow X$ is a covering map, for one can show that it is a proper local homeomorphism. Further, one can find $\eta$, and $\xi$ such that:
\begin{itemize}
    \item $\rho(x, y)\leq 2\eta$ implies $\rho(T(x), T(y)) \geq \lambda\rho(x, y)$
    \item $\rho(x, y)\leq 2\xi$ implies $\rho(T_1^{-1}(x), T_1^{-1}(y)) \leq \frac{1}{\lambda}\rho(x, y)$ where $T_1$ is $T$ restricted to the first slice of the evenly covered neighborhood given by the ball of radius $2\xi$ around $x$.
    \item $B(T(x), \xi) \subset T(B(x, \eta))$
\end{itemize}




\section{Transfer Operator}
Transfer operators provide an operator-theoretic construction of invariant
measures. Begin with
$$m\circ T^{-1} \ll m$$
this gives, 
$$m(T^{-1}(A)) = \int_A \phi dm$$
where $\phi$ is the Radon-Nikodym derivative. If $T$ is a local homeomorphism and the evenly covered neighbourhood $B(x,\xi)$ has slices $U_1,U_2,\ldots,U_n$, then 
$$m\circ \bigg(T\big|_{U_j}\bigg)^{-1} \ll m$$
so, 
$$m(T^{-1}(A)\cap U_j) = m((T|_{U_j})^{-1}(A)) = \int_A \phi_j dm$$
and by the uniqueness of the Radon-Nikodym derivative we have, 
$$\phi = \sum_{j= 1}^{n} \phi_j$$
Consider the map, 
$$\frac{d\mu}{dm} \rightarrow \frac{d(T_*\mu)}{dm}$$
where,
$$\frac{d(T_*\mu)}{dm} = \sum_{j= 1}^{n} \frac{d(\mu\circ(T|_{U_j})^{-1})}{dm} = \sum_{j=1}^{n}\bigg[\frac{d\mu}{dm}\circ(T|_{U_j})^{-1}\bigg]\phi_j$$
Define, 
$$L_m: L^1(m) \rightarrow L^1(m)$$
$$L_m(u)(x):= \sum_{\overline{x}\in T^{-1}(x)}u(\overline{x})\psi(\overline{x})$$
where $\psi \equiv \phi_j\circ T$, and one can check that, 
$$L_m\bigg(\frac{d\mu}{dm}\bigg) = \frac{d(T_*\mu)}{dm}$$
Now, if there exists $h\in L^1(m)$, $h\geq0$, such that $\mu=hm$, then
$$(\mu \text{ is }T\text{-invariant}) \iff (L_m(h) = h)$$
This is an operator-theoretic way of constructing invariant measures. To
construct invariant Gibbs measures for the Ruelle operator $L_{\varphi}$, our
task is twofold: find $m$ such that
\begin{enumerate}
    \item $m\circ T^{-1} \ll m$
    \item $\frac{d(m\circ T^{-1})}{dm} = \text{exp}(\varphi)$
\end{enumerate}
If $L_{\varphi}(h)=h$, then $hm$ is invariant and may also be a Gibbs state. If we see measures as functionals, we have the following,

$$L_{\varphi}^*(\mu)(f) = \mu(L_{\varphi}(f)) = \int_{T(B)}[f\text{exp}(\varphi)]\circ (T|_B)^{-1}d\mu$$
where $\text{supp}(f) \subseteq B$, and for $A\subseteq B$, one can find $\{f_n\}_n \in C(X)$ such that $f_n \rightarrow \chi_A$, which gives, 
$$L_{\varphi}^*(\mu)(A) = \int_{T(A)}\text{exp}(\varphi\circ (T|_A)^{-1})$$
for all $A$ such that $T|_A$ is injective. If $L_{\phi}^*(\mu) = \mu$, then we would have found the measure that satisfies $(1)$ and $(2)$ above. 

\subsection*{Theorem 5.1}
    There exists a measure $m_{\varphi} \in M(X)$ such that,
    $$L_{\phi}^*(\mu) = c\mu$$
    where $c = P(T, \varphi)$.

\textit{Proof.}
    Consider, 
    $$l: M(X) \rightarrow M(X)$$
    $$l(\mu):= \frac{L_{\varphi}^*(\mu)}{L_{\varphi}^*(\mu)(1)}$$
    Now, $M(X)$ is a compact convex set, which implies by the Schauder-Tychnoff theorem: $\exists$ $m$ such that,
    $$l(m) = m$$
    and one can also show that $c = L_{\varphi}^*(\mu)(1) = P(T, \varphi)$.

\subsection*{Theorem 5.2}
    The measure $m$ obtained in Theorem~5.1 is a Gibbs measure for $\varphi$.

\textit{Proof.}
    It can be shown that,
    $$m(B(T^n(x), \xi)) = \int_{T_x^{-n}(B(T^n(x), \xi))}c^n\text{exp}(-S_n\varphi)dm$$
    and following an argument in the previous section we can also show,
    $$1 \geq m(B(T^n(x), \xi)) \geq \alpha$$
    as a result, 
    $$1 \geq \int_{T_x^{-n}(B(T^n(x), \xi))}c^n\text{exp}(-S_n\varphi)dm \geq \alpha$$
    using the pre-bounded distortion lemma, we get two things,
    $$\alpha \leq Cc^n \text{exp}(-S_n\varphi(x)))m(T_x^{-n}(B(T^n(x), \xi)))$$
    $$1 \geq C^{-1}c^n\text{exp}(-S_n\varphi(x))m(T_x^{-n}(B(T^n(x), \xi)))$$
    these put together give,
    $$\frac{\alpha}{C} \leq \frac{m(T_x^{-n}(B(T^n(x), \xi)))}{\text{exp}(S_n\varphi(x) - n\log(c))} \leq C$$

\subsection*{Theorem 5.3}
    Every invariant Gibbs measure is ergodic, and any two Gibbs measures for $\varphi$ are absolutely continuous.

An immediate consequence is that there can be at most one invariant Gibbs
measure. We have found $m_{\varphi}$; it remains to find $h_{\varphi}$ such
that $L_{\varphi}(h_{\varphi})=h_{\varphi}$. Thus
$$h_{\varphi}m_{\varphi}$$
is the required unique Gibbs state for $\varphi$.

\section{Hausdorff Dimension}
A thorough treatment of Hausdorff measures is given in
\cite{rogers1998hausdorff}. For completeness, we give the basic definitions.

\subsection*{Definition 6.1}
    Define, 
    $$H_{\delta}^d(S) := \inf\Bigg\{\sum_{i = 0}^{\infty}(\text{diam}(U_i))^d  \Bigg|  U_i \text{ covers }S, \text{diam}(U_i) < \delta\Bigg\}$$
    and set, 
    $$H^d(S):= \lim_{\delta \rightarrow 0}H_{\delta}^d(S)$$
    The outer measure defined by the function above is the $d$-dimensional Hausdorff measure.


The function
$$\delta \rightarrow H_{\delta}^d(S)$$
is non-increasing. Furthermore, for $s<t$ and every $\delta>0$,
$$H_{\delta}^s(S) \geq \delta^{s-t}H_{\delta}^t(S)$$
\subsection*{Definition 6.2}
    There is a unique number denoted by $HD(S)$ such that,
    \begin{equation*}
        H^d(S) = \begin{cases}
      \infty & \text{if } 0\leq t < HD(S) \\
      0        & \text{if } HD(S) < t <\infty
    \end{cases}
    \end{equation*}
    that number is called the Hausdorff dimension of the set $S$.


Hausdorff measures extend the idea of Lebesgue measure to fractional dimensions. The following covering result is fundamental.

\subsection*{Theorem 6.1}
    Assume that $\mu$ is a Borel probability measure on $\mathbb{R}^n$ and
    that $S$ is a bounded Borel subset of $\mathbb{R}^n$. For $C>0$, the
    following statements hold:
    \begin{enumerate}
        \item for all but countably many $x\in S$, if
        $$\limsup_{r\rightarrow 0} \frac{\mu(B(x, r))}{r^d} \geq C$$
        then $H^d(A) \leq \frac{n}{C}\mu(A)$ for every Borel $A \subset S$.
        \item if for all $x\in S$
        $$\limsup_{r\rightarrow 0} \frac{\mu(B(x, r))}{r^d} \leq C < \infty$$
        then $\mu(A) \leq CH^d(A)$ again for all Borel subsets $A \subset S$.
    \end{enumerate}

\subsection*{Theorem 6.2 — Frostman's Lemma}
    Suppose $\mu$ and $S$ are as above. Then the following hold:
    \begin{enumerate}
        \item If $\mu(S)>0$ and there exists $\theta_1$ such that for every $x\in S$
        $$\liminf_{r\rightarrow 0} \frac{\log(\mu(B(x, r)))}{\log(r)}\geq \theta_1$$
        then $HD(S) \geq \theta_1$
        \item If there exists $\theta_2$ such that for all $x\in S$,
        $$\liminf_{r\rightarrow 0} \frac{\log(\mu(B(x, r)))}{\log(r)} \leq \theta_2$$
        then $HD(S) \leq \theta_2$.
    \end{enumerate}


The preceding density estimates also describe the Hausdorff dimension of a measure; see \cite{falconer2004fractal}. We now introduce conformal repellers and connect them with the ergodic and entropy tools above.

\subsection*{Definition 6.3}
    $(J, V, T)$ is a conformal repeller if for $T:V\rightarrow \mathbb{C}$ (a holomorphic function), one has:
    \begin{enumerate}
        \item $\exists$ $C>0$ and $\alpha > 1$ such that $|(T^n)'(z)| \geq C\alpha$, $\forall$ $z\in J$, and $n\geq 1$.
        \item $T^{-1}(V) \subset \subset V$, $J = \cap_{n\geq 1}T^{-n}(V)$.
        \item $\forall$ $U$, open, $U\cap J \neq \emptyset$, $\exists$ $n>0$ such that $J \subset T^n(U\cap J)$.
    \end{enumerate}

These objects satisfy the standard topological properties summarized below.
\subsection*{Theorem 6.3}
    Given $(J, V, T)$ a conformal repeller, the following statements hold,
    \begin{enumerate}
        \item $T(J) = J$, and $T^{-1}(J) = J$.
        \item $J$ is perfect.
        \item $\exists$ $\beta >0$, such that $\forall$ $x\neq y \in J$, $\exists$ $n$, such that, $|T^n(x) - T^n(y)|\geq \beta$ (i.e. $T$ is expanding on $J$).
        \item $J$ admits a Markov partition $\{J_i\}_{1\leq i\leq k}$.
    \end{enumerate}


\subsection*{Theorem 6.4 — Distortion Lemma}
    Let $m$ be a Gibbs state on $X$, and $\varphi: X\rightarrow \mathbb{R}$ be H\"older-continuous. Assume $P(\varphi) =0$. Then there is a constant $E\geq 1$ such that for all $r$ small enough and all $x\in X$ there exists $n = n(x, r)$ such that,
    $$\frac{\log(E) + S_n\varphi(x)}{-\log(E) - \log(|(T^n)'(x)|)} \leq \frac{\log(m(B(x, r)))}{\log(r)} \leq \frac{-\log(E) + S_n\varphi(x)}{\log(E) - \log(|(T^n)'(x)|)}$$

The distortion estimate yields the volume lemma.
\subsection*{Theorem 6.5 — Volume Lemma}
    Let $m$ be a Gibbs state for a topologically mixing conformal expanding
    repeller on $X$, and let $\varphi$ be a H\"older-continuous real function.
    Then, for almost every $x\in X$, the limit

    $$\lim_{r\rightarrow 0}\frac{\log(m(B(x, r)))}{\log(r)}$$
    exists and is equal to $\frac{h_{\mu}(T)}{\chi_{\mu}(T)}$ (where $\chi_{\mu}(T)$ is the Lyapunov exponent. 

\textit{Proof.}
    Without loss of generality, assume that $P(\varphi) = 0$. We can achieve this by subtracting $P(\varphi)$ from $\varphi$. It can be shown that the Gibbs measure class remains the same up to constant perturbation of the H\"older potential $\varphi$. Now by Birkhoff's Ergodic Theorem, for almost every $x\in X$, we have, 
    $$\lim_{n\rightarrow \infty }\frac{1}{n}S_n\varphi(x) = \int \varphi d\mu$$
    $$\lim_{n\rightarrow \infty} \frac{1}{n}\log|(T^n)'(x)| = \chi_{\mu}(T)$$
    Combining these equalities with the distortion lemma, noting that
    $n\rightarrow\infty$ as $r\rightarrow0$, and using the Variational
    Principle, we obtain

    $$\lim_{r\rightarrow 0}\frac{\log(\mu(B(x, r)))}{\log(r)} = \frac{h_{\mu}(T)}{\chi_{\mu}(T)}$$


This yields the following formula for an invariant ergodic measure $\mu$: 

$$HD(\mu) = \frac{h_{\mu}(T)}{\chi_{\mu}(T)}$$


The Markov partition of a conformal repeller produces a subshift system; see
\cite{zinsmeister2000thermodynamic}. We omit the technical construction and
turn to the result that motivates it.

The function $\phi=-\log|T|'$ is H\"older continuous. Theorem $4.4$ shows
that $\pi(t)=P(t\phi)$ is real analytic and that $\pi(0)$ is the positive
topological entropy of the resulting subshift system. One can further show
that $\pi$ is strictly decreasing, with
$$\lim_{t\rightarrow \infty}\pi(t) = -\infty$$

which implies that there exists a $t>0$ such that $\pi(t) = 0$.
\subsection*{Theorem 6.6 — Bowen} If $t=HD(J)$, then $\pi(t)=0$.
% BLOG-CONTENT-END
\bibliographystyle{amsalpha}
\bibliography{tools/profiles/thermodynamic-formalism/references}
\end{document}
